% comparison-with-baselines.tex -- draft Discussion subsection.
% Consolidates the reflection on the pre-2024 baselines (Section~\ref{sec:baseline})
% against the cohort numbers into four topics. The denominator, median, and window
% points of the earlier eight-point version are merged into the third topic.
% To integrate: paste into the Discussion of paper.tex, after "External verification".

\subsection{What the baseline comparison does and does not show}
\label{sec:baseline-comparison}

\textbf{The cohort sits outside every historical envelope, including the tails.}
The baseline table (\autoref{tab:baseline}) aggregates studies with different
populations and selection rules, so the honest comparison for a curated cohort is
against the most extreme rows, not the averages.
Against GitClear's senior developers at large technology companies (4.28 commits
per day), its 99th percentile (up to ten times the median developer, roughly 17
commits per working day), and its hard ceiling of about 55 commits per day, the
cohort still stands outside: the least active member makes 7 commits per active
day, the median member 31, and the two most active 188 and 304
(\autoref{tab:rq4-breadth}).
That ceiling is itself a finding about the baselines: GitClear excludes accounts
above roughly 55 commits per day ``presumably because such output is considered
implausible for a human and characteristic of automation such as bots.''
Two of the nine developers sustain 188 and 304 commits per active day---verified
humans directing agents, not bots.
The baseline numbers encode an assumption about human limits, and the cohort
falsifies it.
The defensible claim is therefore not ``$N$ times the baseline'' but ``outside
every envelope in the literature, including the prolific rows, the 99th
percentile, and the plausibility ceiling.''

\textbf{The baselines no longer measure the same thing, and the newest ones are
already AI-era.}
Once agents commit autonomously, commit frequency, work schedules, and repository
switching stop measuring developer behavior and start measuring the output of a
developer--agent system: commits land during the assumed rest hours (RQ1), and
the repository-switching counts of RQ3 measure the recorded activity, not the
developer's presence at a keyboard.
The metrics keep their names but change their constructs, so the comparison is
system against system, not developer against developer.
At the same time, the baseline end of the comparison is contaminated upward: the
baselines ``mostly cover data collected before 2024, with some studies including
data from 2024 partially,'' but GitClear's window runs to 2024, and the Stray
et al.\ and Karyotakis et al.\ entries are 2026 publications measured in the
Copilot era.
The measured gap is therefore a lower bound on the true pre-agent gap.
The consequence for research practice is that activity baselines need to be
re-derived for the agent era; until then, cross-era comparisons must be framed
as system output against historical system output.

\textbf{Comparing fairly: denominators, medians, and windows.}
Three methodological points discipline the comparison.
First, denominators: the baselines mix commits per day~\cite{kolassa2013empirical},
per working day~\cite{gitclear2024git}, per week over 105
weeks~\cite{stray2026copilot}, per active hour~\cite{zhang2024workrhythms}, and
per project-year~\cite{milewicz2019characterizing}.
The cohort's per-active-day rates survive any conversion for the median and top
members (teamchong: 31 per active day, about 30 per calendar day), but not for the
least active member: philipp-spiess is 7 per active day---above the prolific
rows---yet about 3--4 per calendar day, inside them.
The denominator must be stated per row, and it matters at the cohort's low end.
Second, medians: both sides are long-tailed (Kolassa's model reports 16 LoC median
against 465.7 mean~\cite{kolassa2013model}; the cohort's largest commit changes
6.9 million lines), and the tails are different things---machine state and
vendored artifacts on our side, human mega-commits on theirs.
Mean-based comparison is tail-dominated and unfair; the median comparison
(27--88 lines against roughly 20--24 LoC~\cite{riehle2012developer,sadowski2018modern})
is the defensible one, and it is what makes the volume gap robust to commit
granularity.
Third, windows: the baselines average over years, which smooths bursts, while the
cohort window is eight months and burst-collapse is modal (teamchong: 3{,}147
commits in one month, then a decline).
Per-active-day rates are the right cohort statistic, and yearly normalization
would smooth the cohort's bursts downward---the window does not inflate the gap.

\textbf{What the comparison proves, and what it does not.}
The comparison establishes a plausibility bound: the cohort operates outside every
historical envelope we could find.
It does not establish causation: the populations differ (curated by us;
longevity-conditioned in Vasilescu et al.; plausibility-filtered in GitClear),
the eras differ, and the ecosystems differ.
The missing baseline is the important one: there is no measurement of developers
\emph{with} agents at ordinary volumes.
Only that baseline would turn ``the envelope has moved'' into ``agents move it by
$N$.''
Until such a baseline exists, the honest statement of the productivity finding is
the bound, not a multiplier.

The actionable consequences of these four points, for practitioners and for
researchers, are stated in the implications subsections that follow.
